\chapter{Теория}

\section{nslookup Ubuntu}

Необходимо выполнить команду

\begin{lstlisting}
    sudo nslookup <domain_name> # например sudo nslookup yandex.ru
\end{lstlisting}

\begin{figure}[H] 
    \centering
    \includegraphics[width=1.0\textwidth]{nsl.png}
    \caption{Вывод команды nslookup}
\end{figure}

Поле «Server» указывает IP-адрес DNS-сервера, к которому
обратился компьютер (ПК), а затем выводится информация об IP-адресе домена yandex.ru
(выделено синим).
В данном результате присутствует поле «Non - Authoritative Answer» (неавторитетный
ответ, резервный). Авторитетным (основным) считается ответ от DNS-сервера, на котором
содержится полная информация о зоне домена. Во многих случаях DNS-серверы не являются
авторитетными и хранят в памяти результаты предыдущих запросов, на которые были
получены авторитетные ответы. Когда такой сервер получает запрос, он ищет в памяти
соответствующую информацию и, при наличии, отправляет ее в качестве неавторитетного
ответа.

\subsection{Ключ -query=mx}

\begin{figure}[H] 
    \centering
    \includegraphics[width=1.0\textwidth]{nsq.png}
    \caption{Применение флага -query=mx}
\end{figure}

Запись MX (Mail eXchange, обмен почтой) хранит соответствие доменного имени
почтовому серверу этого домена. В примере вывод команды указал почтовый сервер домена,
через который должна отправляться вся электронная почта на адреса «@yandex.ru». Для
3
домена yandex.ru есть одна запись МХ. Число рядом с именем сервера (10) означает его
приоритет.


\subsection{Ключ -query=ns}

\begin{figure}[H] 
    \centering
    \includegraphics[width=1.0\textwidth]{ns.png}
    \caption{Применение флага -query=ns}
\end{figure}

Запись NS (Name Server, сервер имен) содержит информацию о соответствии между
доменным именем и авторитетным (основным) DNS-сервером для данного домена.

\subsection{Ключ -query=soa}

\begin{figure}[H] 
    \centering
    \includegraphics[width=1.0\textwidth]{soa.png}
    \caption{Применение флага -query=soa}
\end{figure}

Запись SOA (Start of Authority, начальная запись зоны) содержит информацию о зоне
домена, адрес его администратора, серийный номер и т.д. На примере:
\begin{enumerate}
    \item origin – имя первичного сервера зоны;
    \item  mail addr – адрес администратора домена (первая точка в адресе заменяет символ @ в описании зоны, так как он имеет свое значение);
    \item serial – серийный номер файла зоны, используется для учета изменений. Может быть любое целое число, но стандартный формат — «ГГГГММДДНН», то есть сначала указывается дата, а НН (в данном случае 21) увеличивается в случае нескольких обновлений в день;
    \item refresh – период времени (в секундах), через который вторичный DNS-сервер отправит запрос первичному, чтобы проверить, поменялся ли серийный номер;
    \item retry – указывает интервал для повторного соединения с первичным DNSсервером, если он по каким-то причинам не смог ответить на запрос;
    \item expire – указывает время хранения кэша вторичным DNS-сервером, по истечении которого он будет считаться устаревшим;
    \item minimum – минимальное время хранения кэша вторичным DNS-сервером до повторного запроса.
\end{enumerate}

\subsection{Ключ -query=any}

\begin{figure}[H] 
    \centering
    \includegraphics[width=1.0\textwidth]{any.png}
    \caption{Применение флага -query=any}
\end{figure}

С помощью данного ключа мы можно просмотреть все записи DNS, которые есть для
заданного имени. Как мы видим, это почтовые сервера, доверенные сервера, IP-адреса
различных серверов связанных с этим доменным именем (в том числе IPv-6 адреса).

\subsection{Обратный поиск}

\begin{figure}[H] 
    \centering
    \includegraphics[width=1.0\textwidth]{r.png} 
    \caption{Обратный поиск}
\end{figure}

Если вместо имени домена указать его IP адрес, то будет выполнен обратный поиск. Необходимо отметить, в данном результате отсутствует фраза "Non-authoritative answer"
в связи с тем, что DNS-сервер ns1.yandex.ru обладает полной информацией о зоне для
домена yandex.ru.

\subsection{Интерактивный режим}
Для входа в интерактивный режим запустите команду nslookup без опций.


\endinput